% --- Archon physics formalization source begin ---
% archon:physics
% archon:covers PhyXMiniProblems/problem_phyx_mini_0194.lean
% archon:source-report reports/phyx_mini/problem_phyx_mini_0194.source.json

\chapter{Physics problem phyx\_mini\_0194}
\label{ch:PhyXMiniProblems_problem_phyx_mini_0194}

\paragraph{Problem source.}
\#\# Physical scenario

A listener L is at rest and the siren is moving away from L at 30 \textbackslash{}, \textbackslash{}mathrm\{m/s\}.

\#\# Question

What frequency does the listener hear?

\#\# Answer choices

- A: A: 267Hz

- B: B: 276Hz

- C: C: 283Hz

- D: D: 296Hz

\#\# Figure caption (auxiliary; use the image as primary evidence)

The image depicts a scenario illustrating the Doppler effect. It shows a police car and a listener.

1. **Listener at Left:**
   - The listener is depicted as a simple stick figure on the left, labeled "L."
   - Above this figure, there is text saying "Listener at rest" with "v\_L = 0" indicating the listener is stationary.
   - It is also asked "f\_L = ?," indicating the frequency perceived by the listener is unknown.

2. **Police Car:**
   - The police car is depicted on the right, labeled "S."
   - An arrow above the car points to the right, indicating the direction of its movement.
   - Above the car, there's the text "Police car" and "v\_S = 30 m/s," indicating the speed of the car.

3. **Arrow and Text:**
   - There’s an arrow labeled "L to S" pointing from the listener to the car, indicating the direction from the listener to the source.

Overall, the image visualizes a typical Doppler effect scenario where a source of sound (the police car) moves towards a stationary observer (the listener).

\#\# Dataset metadata

Wave Properties; Physical Model Grounding Reasoning

\paragraph{Recorded answer/context.}
B: B: 276Hz

\paragraph{Figure/image path.}
/root/proposal\_for\_physic/science-mango/phyx\_mini\_run/phyx\_data/test\_image/194.png

\paragraph{Formalization target.}
create a compiling Lean file with sorry bodies at `PhyXMiniProblems/problem\_phyx\_mini\_0194.lean`. The Lean declarations must preserve the physical quantities, dimensions or dimensional roles, figure labels, governing-law hypotheses, and final relation expressed by this problem.
Use Mathlib/Physlib names found through LeanExplore where available. If a physics API is missing, introduce faithful local abstractions rather than scalar placeholder aliases.

\begin{theorem}[Physics formalization target]
\label{thm:physics:phyx_mini_0194:target}
\lean{PhyXMiniProblems.ProblemPhyXMini0194.recedingSiren_matches_recordedAnswerB}
\uses{def:physics:phyx-mini-0194:phyxminiproblems-problemphyxmini0194-recedingsirensetup, def:physics:phyx-mini-0194:phyxminiproblems-problemphyxmini0194-matchesproblemandprimaryfigure, def:physics:phyx-mini-0194:phyxminiproblems-problemphyxmini0194-usesstandardstillair, def:physics:phyx-mini-0194:phyxminiproblems-problemphyxmini0194-usesnominalsirenemission, def:physics:phyx-mini-0194:phyxminiproblems-problemphyxmini0194-hasphysicaldopplerparameters, def:physics:phyx-mini-0194:phyxminiproblems-problemphyxmini0194-satisfiesleftwardacousticdopplerlaw, def:physics:phyx-mini-0194:phyxminiproblems-problemphyxmini0194-recordedanswerchoice, def:physics:phyx-mini-0194:phyxminiproblems-problemphyxmini0194-matchesanswerchoice}
The assigned autoformalize agent should translate this physics problem into Lean declarations in the covered file, with theorem and lemma proof bodies written as `by sorry`.
\end{theorem}
\begin{proof}
This is an autoformalization task, not a proof task. Produce statements that can later be proved by the physics prover without weakening the physical model.
\end{proof}

% --- Archon declaration topology begin ---
\section{Declaration topology}
\begin{definition}[Acoustic Frequency]
  \label{def:physics:phyx-mini-0194:phyxminiproblems-problemphyxmini0194-acousticfrequency}
  \lean{PhyXMiniProblems.ProblemPhyXMini0194.AcousticFrequency}
  A nonnegative physical acoustic frequency, with inverse-time dimension
\end{definition}
\begin{proof}
  This is the declared model definition of Acoustic Frequency.
\end{proof}
\begin{definition}[Axial Position]
  \label{def:physics:phyx-mini-0194:phyxminiproblems-problemphyxmini0194-axialposition}
  \lean{PhyXMiniProblems.ProblemPhyXMini0194.AxialPosition}
  A signed physical coordinate along the horizontal L-to-S axis
\end{definition}
\begin{proof}
  This is the declared model definition of Axial Position.
\end{proof}
\begin{definition}[Axial Velocity]
  \label{def:physics:phyx-mini-0194:phyxminiproblems-problemphyxmini0194-axialvelocity}
  \lean{PhyXMiniProblems.ProblemPhyXMini0194.AxialVelocity}
  A signed physical velocity along the horizontal axis. Positive velocity points from listener L toward source S in the primary figure
\end{definition}
\begin{proof}
  This is the declared model definition of Axial Velocity.
\end{proof}
\begin{definition}[frequency Readout]
  \label{def:physics:phyx-mini-0194:phyxminiproblems-problemphyxmini0194-frequencyreadout}
  \lean{PhyXMiniProblems.ProblemPhyXMini0194.frequencyReadout}
  \uses{def:physics:phyx-mini-0194:phyxminiproblems-problemphyxmini0194-acousticfrequency}
  Read a physical acoustic frequency in inverse units of the selected time unit
\end{definition}
\begin{proof}
  This is the declared model definition of frequency Readout.
\end{proof}
\begin{definition}[frequency In Hertz]
  \label{def:physics:phyx-mini-0194:phyxminiproblems-problemphyxmini0194-frequencyinhertz}
  \lean{PhyXMiniProblems.ProblemPhyXMini0194.frequencyInHertz}
  \uses{def:physics:phyx-mini-0194:phyxminiproblems-problemphyxmini0194-acousticfrequency, def:physics:phyx-mini-0194:phyxminiproblems-problemphyxmini0194-frequencyreadout}
  Hertz readout of a physical acoustic frequency
\end{definition}
\begin{proof}
  This is the declared model definition of frequency In Hertz.
\end{proof}
\begin{definition}[position Readout]
  \label{def:physics:phyx-mini-0194:phyxminiproblems-problemphyxmini0194-positionreadout}
  \lean{PhyXMiniProblems.ProblemPhyXMini0194.positionReadout}
  \uses{def:physics:phyx-mini-0194:phyxminiproblems-problemphyxmini0194-axialposition}
  Read a signed axial position in the selected length unit
\end{definition}
\begin{proof}
  This is the declared model definition of position Readout.
\end{proof}
\begin{definition}[position In Meters]
  \label{def:physics:phyx-mini-0194:phyxminiproblems-problemphyxmini0194-positioninmeters}
  \lean{PhyXMiniProblems.ProblemPhyXMini0194.positionInMeters}
  \uses{def:physics:phyx-mini-0194:phyxminiproblems-problemphyxmini0194-axialposition, def:physics:phyx-mini-0194:phyxminiproblems-problemphyxmini0194-positionreadout}
  Metre readout of an axial position
\end{definition}
\begin{proof}
  This is the declared model definition of position In Meters.
\end{proof}
\begin{definition}[axial Velocity Readout]
  \label{def:physics:phyx-mini-0194:phyxminiproblems-problemphyxmini0194-axialvelocityreadout}
  \lean{PhyXMiniProblems.ProblemPhyXMini0194.axialVelocityReadout}
  \uses{def:physics:phyx-mini-0194:phyxminiproblems-problemphyxmini0194-axialvelocity}
  Read a signed axial velocity in selected length and time units
\end{definition}
\begin{proof}
  This is the declared model definition of axial Velocity Readout.
\end{proof}
\begin{definition}[axial Velocity In Meters Per Second]
  \label{def:physics:phyx-mini-0194:phyxminiproblems-problemphyxmini0194-axialvelocityinmeterspersecond}
  \lean{PhyXMiniProblems.ProblemPhyXMini0194.axialVelocityInMetersPerSecond}
  \uses{def:physics:phyx-mini-0194:phyxminiproblems-problemphyxmini0194-axialvelocity, def:physics:phyx-mini-0194:phyxminiproblems-problemphyxmini0194-axialvelocityreadout}
  Metres-per-second readout of a signed axial velocity
\end{definition}
\begin{proof}
  This is the declared model definition of axial Velocity In Meters Per Second.
\end{proof}
\begin{definition}[speed In Meters Per Second]
  \label{def:physics:phyx-mini-0194:phyxminiproblems-problemphyxmini0194-speedinmeterspersecond}
  \lean{PhyXMiniProblems.ProblemPhyXMini0194.speedInMetersPerSecond}
  Read Physlib's nonnegative dimensionful speed in metres per second
\end{definition}
\begin{proof}
  This is the declared model definition of speed In Meters Per Second.
\end{proof}
\begin{definition}[Figure Point]
  \label{def:physics:phyx-mini-0194:phyxminiproblems-problemphyxmini0194-figurepoint}
  \lean{PhyXMiniProblems.ProblemPhyXMini0194.FigurePoint}
  The two labelled locations in the primary figure
\end{definition}
\begin{proof}
  This is the declared model definition of Figure Point.
\end{proof}
\begin{definition}[Sound Source Kind]
  \label{def:physics:phyx-mini-0194:phyxminiproblems-problemphyxmini0194-soundsourcekind}
  \lean{PhyXMiniProblems.ProblemPhyXMini0194.SoundSourceKind}
  The physical kind of sound source drawn at S
\end{definition}
\begin{proof}
  This is the declared model definition of Sound Source Kind.
\end{proof}
\begin{definition}[Axial Direction]
  \label{def:physics:phyx-mini-0194:phyxminiproblems-problemphyxmini0194-axialdirection}
  \lean{PhyXMiniProblems.ProblemPhyXMini0194.AxialDirection}
  The two orientations along the one-dimensional figure axis
\end{definition}
\begin{proof}
  This is the declared model definition of Axial Direction.
\end{proof}
\begin{definition}[Receding Siren Setup]
  \label{def:physics:phyx-mini-0194:phyxminiproblems-problemphyxmini0194-recedingsirensetup}
  \lean{PhyXMiniProblems.ProblemPhyXMini0194.RecedingSirenSetup}
  \uses{def:physics:phyx-mini-0194:phyxminiproblems-problemphyxmini0194-acousticfrequency, def:physics:phyx-mini-0194:phyxminiproblems-problemphyxmini0194-axialposition, def:physics:phyx-mini-0194:phyxminiproblems-problemphyxmini0194-axialvelocity, def:physics:phyx-mini-0194:phyxminiproblems-problemphyxmini0194-figurepoint, def:physics:phyx-mini-0194:phyxminiproblems-problemphyxmini0194-soundsourcekind, def:physics:phyx-mini-0194:phyxminiproblems-problemphyxmini0194-axialdirection}
  Independent physical quantities and qualitative labels for the Doppler experiment. In particular, heardFrequencyAtL is an unknown physical frequency and is not assigned a numerical value by this structure
\end{definition}
\begin{proof}
  This is the declared model definition of Receding Siren Setup.
\end{proof}
\begin{definition}[Matches Problem And Primary Figure]
  \label{def:physics:phyx-mini-0194:phyxminiproblems-problemphyxmini0194-matchesproblemandprimaryfigure}
  \lean{PhyXMiniProblems.ProblemPhyXMini0194.MatchesProblemAndPrimaryFigure}
  \uses{def:physics:phyx-mini-0194:phyxminiproblems-problemphyxmini0194-positioninmeters, def:physics:phyx-mini-0194:phyxminiproblems-problemphyxmini0194-axialvelocityinmeterspersecond, def:physics:phyx-mini-0194:phyxminiproblems-problemphyxmini0194-recedingsirensetup}
  Problem-statement and primary-figure data. The listener L is left of S and at rest; source S is a police-car siren moving at 30 m/s in the L-to-S direction, hence away from L. The sound that reaches L travels from S toward L. No emitted-frequency calibration, sound-speed value, or requested heard frequency occurs in this predicate
\end{definition}
\begin{proof}
  This is the declared model definition of Matches Problem And Primary Figure.
\end{proof}
\begin{definition}[Uses Standard Still Air]
  \label{def:physics:phyx-mini-0194:phyxminiproblems-problemphyxmini0194-usesstandardstillair}
  \lean{PhyXMiniProblems.ProblemPhyXMini0194.UsesStandardStillAir}
  \uses{def:physics:phyx-mini-0194:phyxminiproblems-problemphyxmini0194-axialvelocityinmeterspersecond, def:physics:phyx-mini-0194:phyxminiproblems-problemphyxmini0194-speedinmeterspersecond, def:physics:phyx-mini-0194:phyxminiproblems-problemphyxmini0194-recedingsirensetup}
  Standard still-air calibration used for the numerical multiple-choice evaluation. It is kept separate because neither still air nor 343 m/s is printed in the supplied problem or bitmap
\end{definition}
\begin{proof}
  This is the declared model definition of Uses Standard Still Air.
\end{proof}
\begin{definition}[Uses Nominal Siren Emission]
  \label{def:physics:phyx-mini-0194:phyxminiproblems-problemphyxmini0194-usesnominalsirenemission}
  \lean{PhyXMiniProblems.ProblemPhyXMini0194.UsesNominalSirenEmission}
  \uses{def:physics:phyx-mini-0194:phyxminiproblems-problemphyxmini0194-frequencyinhertz, def:physics:phyx-mini-0194:phyxminiproblems-problemphyxmini0194-recedingsirensetup}
  Nominal source-frequency calibration required by the supplied answer table. The source and bitmap omit this datum, so it is deliberately separated from MatchesProblemAndPrimaryFigure
\end{definition}
\begin{proof}
  This is the declared model definition of Uses Nominal Siren Emission.
\end{proof}
\begin{definition}[Has Physical Doppler Parameters]
  \label{def:physics:phyx-mini-0194:phyxminiproblems-problemphyxmini0194-hasphysicaldopplerparameters}
  \lean{PhyXMiniProblems.ProblemPhyXMini0194.HasPhysicalDopplerParameters}
  \uses{def:physics:phyx-mini-0194:phyxminiproblems-problemphyxmini0194-frequencyinhertz, def:physics:phyx-mini-0194:phyxminiproblems-problemphyxmini0194-axialvelocityinmeterspersecond, def:physics:phyx-mini-0194:phyxminiproblems-problemphyxmini0194-speedinmeterspersecond, def:physics:phyx-mini-0194:phyxminiproblems-problemphyxmini0194-recedingsirensetup}
  Positivity and subsonic conditions for the classical acoustic model
\end{definition}
\begin{proof}
  This is the declared model definition of Has Physical Doppler Parameters.
\end{proof}
\begin{definition}[Satisfies Leftward Acoustic Doppler Law]
  \label{def:physics:phyx-mini-0194:phyxminiproblems-problemphyxmini0194-satisfiesleftwardacousticdopplerlaw}
  \lean{PhyXMiniProblems.ProblemPhyXMini0194.SatisfiesLeftwardAcousticDopplerLaw}
  \uses{def:physics:phyx-mini-0194:phyxminiproblems-problemphyxmini0194-frequencyinhertz, def:physics:phyx-mini-0194:phyxminiproblems-problemphyxmini0194-axialvelocityinmeterspersecond, def:physics:phyx-mini-0194:phyxminiproblems-problemphyxmini0194-speedinmeterspersecond, def:physics:phyx-mini-0194:phyxminiproblems-problemphyxmini0194-recedingsirensetup}
  The one-dimensional classical Doppler law for sound travelling leftward from S to L. Signed listener and source velocities are measured relative to the air. This governing law relates independent physical quantities and does not mention 276 Hz or any answer label
\end{definition}
\begin{proof}
  This is the declared model definition of Satisfies Leftward Acoustic Doppler Law.
\end{proof}
\begin{definition}[Answer Choice]
  \label{def:physics:phyx-mini-0194:phyxminiproblems-problemphyxmini0194-answerchoice}
  \lean{PhyXMiniProblems.ProblemPhyXMini0194.AnswerChoice}
  Labels of the four answers printed with the problem
\end{definition}
\begin{proof}
  This is the declared model definition of Answer Choice.
\end{proof}
\begin{definition}[hertz]
  \label{def:physics:phyx-mini-0194:phyxminiproblems-problemphyxmini0194-answerchoice-hertz}
  \lean{PhyXMiniProblems.ProblemPhyXMini0194.AnswerChoice.hertz}
  \uses{def:physics:phyx-mini-0194:phyxminiproblems-problemphyxmini0194-answerchoice}
  Whole-hertz frequency printed beside each answer label
\end{definition}
\begin{proof}
  This is the declared model definition of hertz.
\end{proof}
\begin{definition}[recorded Answer Choice]
  \label{def:physics:phyx-mini-0194:phyxminiproblems-problemphyxmini0194-recordedanswerchoice}
  \lean{PhyXMiniProblems.ProblemPhyXMini0194.recordedAnswerChoice}
  \uses{def:physics:phyx-mini-0194:phyxminiproblems-problemphyxmini0194-answerchoice}
  The recorded answer label in the supplied dataset metadata
\end{definition}
\begin{proof}
  This is the declared model definition of recorded Answer Choice.
\end{proof}
\begin{definition}[Matches Answer Choice]
  \label{def:physics:phyx-mini-0194:phyxminiproblems-problemphyxmini0194-matchesanswerchoice}
  \lean{PhyXMiniProblems.ProblemPhyXMini0194.MatchesAnswerChoice}
  \uses{def:physics:phyx-mini-0194:phyxminiproblems-problemphyxmini0194-acousticfrequency, def:physics:phyx-mini-0194:phyxminiproblems-problemphyxmini0194-frequencyinhertz, def:physics:phyx-mini-0194:phyxminiproblems-problemphyxmini0194-answerchoice}
  Agreement with a displayed whole-hertz answer to half a hertz, representing rounding to the nearest integer rather than an exact-integer physical value
\end{definition}
\begin{proof}
  This is the declared model definition of Matches Answer Choice.
\end{proof}
\begin{lemma}[heard Frequency In Hertz eq exact Model Value]
  \label{lem:physics:phyx-mini-0194:phyxminiproblems-problemphyxmini0194-heardfrequencyinhertz-eq-exactmodelvalue}
  \lean{PhyXMiniProblems.ProblemPhyXMini0194.heardFrequencyInHertz_eq_exactModelValue}
  \uses{def:physics:phyx-mini-0194:phyxminiproblems-problemphyxmini0194-frequencyinhertz, def:physics:phyx-mini-0194:phyxminiproblems-problemphyxmini0194-recedingsirensetup, def:physics:phyx-mini-0194:phyxminiproblems-problemphyxmini0194-matchesproblemandprimaryfigure, def:physics:phyx-mini-0194:phyxminiproblems-problemphyxmini0194-usesstandardstillair, def:physics:phyx-mini-0194:phyxminiproblems-problemphyxmini0194-usesnominalsirenemission, def:physics:phyx-mini-0194:phyxminiproblems-problemphyxmini0194-satisfiesleftwardacousticdopplerlaw}
  Under the stated figure data and the two explicit auxiliary calibrations, the classical Doppler model predicts the exact SI readout 102900 / 373 Hz
\end{lemma}
\begin{proof}
  The conclusion follows from the governing relations and figure readouts named in the statement of heard Frequency In Hertz eq exact Model Value.
\end{proof}
% --- Archon declaration topology end ---
% --- Archon physics formalization source end ---
